\usepackage{environ} 
\usepackage{lastpage} 
\usepackage{fancyhdr} 
\usepackage{tcolorbox}
\tcbuselibrary{skins,xparse,breakable,raster,listings,listingsutf8,documentation,many}
\usepackage{xcolor}
\usepackage{amsmath}
\usepackage{graphicx}

\def\vec{\overrightarrow}
\newcommand{\hoac}[1]{\left[\begin{aligned}#1\end{aligned}\right.}
\newcommand{\heva}[1]{\left\{\begin{aligned}#1\end{aligned}\right.}
\newcommand{\hetde}{\centerline{\rule[0.5ex]{2cm}{1pt} HẾT \rule[0.5ex]{2cm}{1pt}}}

\newcommand{\tentruong}{THPT Nguyễn Hữu Cảnh}
\newcommand{\tengv}{GV Nguyễn Văn Sang}
\newcommand{\tenkythi}{Ôn Kiểm Tra GHKII 2025}
\newcommand{\tenmonthi}{Môn: Toán - Khối 12}
\newcommand{\thoigian}{90}
\newcommand{\made}{101}

\newcommand{\tieude}[3]{
	\noindent
	% Trái
	\begin{minipage}[b]{7cm}
		\centerline{\textbf{\fontsize{13}{0}\selectfont \tentruong}}
		\centerline{\textbf{\fontsize{13}{0}\selectfont \tengv}}
		\centerline{(\textit{Đề thi có #1\ trang})}
		\centerline{\textit{ }}
	\end{minipage}\hspace{1.5cm}
	% Phải
	\begin{minipage}[b]{9cm}
		\centerline{\textbf{\fontsize{13}{0}\selectfont \tenkythi}}
		\centerline{\textbf{\fontsize{13}{0}\selectfont \tenmonthi}}
		\centerline{\textit{\fontsize{12}{0}\selectfont Thời gian làm bài \thoigian \ phút }}
		% Đã sửa lỗi Overfull \hbox bằng cách ngắt dòng
		\centerline{\textit{\fontsize{12}{0}\selectfont (Đề có 11 câu trắc nghiệm, 5 câu đúng sai, 6 câu điền đáp án)}}
	\end{minipage}
\vspace{0.5em}

	\begin{minipage}[b]{13cm}
		\vspace*{.75cm}
		\textbf{Họ và tên thí sinh: }{\tiny\dotfill}\\
		\textbf{Số Báo Danh: }{\tiny\dotfill}
	\end{minipage}
	\begin{minipage}[b]{8cm}
		\hspace*{1cm}\fbox{\bf Mã đề thi #3}
	\end{minipage}
}
% Định nghĩa lại chantrang để thiết lập footer cho fancyhdr
\newcommand{\chantrang}[2]{
    \fancyhf{} % Xóa tất cả các header và footer mặc định
    \rfoot{Trang \thepage/#1 $-$ Mã đề #2} % Đặt nội dung vào chân trang bên phải
    \lfoot{} % Để trống chân trang bên trái
    \cfoot{} % Để trống chân trang ở giữa
}
\renewcommand{\headrulewidth}{0pt}
% --- Các định nghĩa ---
\newcommand{\caulc}{
	\begin{flushleft}
{\color{red}\sffamily\bfseries Phần I. Câu hỏi trắc nghiệm nhiều phương án lựa chọn}\\
\textit{Thí sinh trả lời từ câu 1 đến câu 12. Mỗi câu hỏi thí sinh chỉ chọn một phương án.}
	\end{flushleft}
\setcounter{ex}{0}


	\setcounter{ex}{0}
}
\newcommand{\cauds}{
	\begin{flushleft}
{\color{red}\sffamily\bfseries Phần II. Câu hỏi trắc nghiệm đúng sai}\\
\textit{Thí sinh trả lời từ câu 1 đến câu 4. Trong mỗi ý a), b), c), d) ở mỗi câu thí sinh chọn đúng hoặc sai.} 
	\end{flushleft}
	\setcounter{ex}{0}
}
\newcommand{\caukq}{
	\begin{flushleft}
{\color{red}\sffamily\bfseries Phần III. Câu hỏi trắc nghiệm trả lời ngắn}\\
\textit{Thí sinh trả lời từ câu 1 đến câu 6.} 	\end{flushleft}
	\setcounter{ex}{0}
} 

% === CƠ CHẾ ẨN/HIỆN LỜI GIẢI ===
\newif\ifhienloigiai

\newcommand{\hienloigiai}{
    \hienloigiaitrue
    \makeatletter
    \def\colorEX{\color{red}}
    \renewcommand{\TrueEX}{\stepcounter{dapan}\squareEX{\bf\Alph{dapan}}\ignorespaces}
    \makeatother
    \typeout{>>> CHE DO DAP AN DA DUOC KICH HOAT.}
}

\newcommand{\anloigiai}{
    \hienloigiaifalse
    \makeatletter
    \def\colorEX{}
    \renewcommand{\TrueEX}{\FalseEX}
    \makeatother
    \typeout{>>> CHE DO DE THI DA DUOC KICH HOAT.}
}

\RenewEnviron{loigiai}{%
    \ifhienloigiai
        \par\noindent\textbf{\loigiaiEX}
        \BODY
    \fi
    \makeatletter
    \xdef\keyEX{}%
    \setcounter{numTrueSol}{0}%
    \makeatother
}
